\documentclass[a4paper,12pt]{article}
\usepackage[utf8]{inputenc}
\usepackage[ngerman]{babel}
\usepackage{amsmath}
\usepackage{geometry}
\geometry{a4paper, top=15mm, left=15mm, right=15mm, bottom=15mm, 
	headsep=10mm, footskip=12mm}
\usepackage{parskip}
\usepackage{circuitikz}
\usepackage{graphicx}
\usepackage{tikz}
\usepackage{href-ul}
\hypersetup{ 
	colorlinks=true, 
	linkcolor=blue, 
	urlcolor=blue}

\begin{document} 
	
	\begin{tikzpicture}(10,3)
	\draw[ultra thick](2,0) --(10,0) -- (10,1.3) --(2,1.3) -- (2,0);
	\draw[fill=black](2,0)-- (0,.6) -- (2,.6) -- (2,0);
	\node at (6,.6) {\large Solution Sheet of \href{https://www.okuyakl.de}{www.okuyakl.de}};
    \end{tikzpicture}
	
	\section*{Basic Knowledge – Solutions}
	
	\begin{itemize}
		\item \textbf{A closed circuit} is necessary for a lamp to light up.
		
		\item The current flows from the \textbf{positive} pole of the battery to the \textbf{negative} pole (technical current direction).
		
		\item In an electrical circuit, you need at least:
		\begin{itemize}
			\item( x ) a battery or voltage source
			\item( x ) a load (e.g., lamp)
			\item( x ) cables or wires
		\end{itemize}
		
		\item Symbols:
		
		\textbf{Battery:}
		\begin{circuitikz}
			\draw (0,0) to[battery] (2,0);
		\end{circuitikz}
		
		\textbf{Lamp:}
		\begin{circuitikz}
			\draw (0,0) to[lamp] (2,0);
		\end{circuitikz}
		
		\textbf{Switch (open):}
		\begin{circuitikz}
			\draw (0,0) to[switch] (2,0);
		\end{circuitikz}
	\end{itemize}
	
	\hrule
	
	\section*{Understanding Electrical Circuits }
	
	\begin{itemize}
		\item a) Why does the circuit have to be "closed"? - So that the electrons can flow unhindered. The voltage source acts like an electron pump, creating a shortage of electrons at the positive pole and an excess of electrons at the negative pole.
		
		\vspace{1em}
		\underline{\hspace{12cm}}
		
		\item b) What are the possible circuit options?
		
		\begin{minipage}{0.4 \textwidth}
			
			Series circuit:
			
			\begin{circuitikz}[american voltages]
				% battery on the left, two lamps in series at the top, return at the bottom
				\draw
				(0,0) to[battery1,l=$V$] (0,3) % voltage source
				-- (1.5,3) to[lamp,l=$L_1$] (2.5,3) % first bulb
				to[lamp,l=$L_2$] (6,3) % second bulb
				-- (6,0) -- (0,0); % Return line to the negative pole
			\end{circuitikz}
			
		\end{minipage}
		\begin{minipage}{0.4 \textwidth}
			
			Parallel connection:
			
			\begin{circuitikz}[american voltages]
				% Battery left
				\draw
				(0,0) to[battery1,l=$V$] (0,4)
				-- (4,4) % Upper line
				(0,0) -- (4,0); % Lower line
				
				% Two lamps in parallel (vertical)
				\draw
				(2,4) to[lamp,l=$L_1$] (2,0);
				\draw
				(4,4) to[lamp,l=$L_2$] (4,0);
			\end{circuitikz}
		\end{minipage}
		
		\newpage
		\item c) Check:
		In a series circuit of two lamps...
		
		\begin{itemize}
			\item( ) both lamps can be turned on and off independently.
			\item(x) both lamps go out when one burns out.
			\item(x) the same current flows through all lamps, even if they are different.
			\item( ) it gets brighter if you add more lamps.
			\item(x) both lamps must have the same operating current in amperes.
			\item( ) both lamps must have the same operating voltage in volts.
			\item( ) both lamps must have the same power in watts.
			\item(x) the sum of the operating voltages of the lamps must be equal to the voltage of the power source.
		\end{itemize}
		
		\item d) Check:
		In a parallel circuit of two lamps...
		
		\begin{itemize}
			\item(x) both lamps can be turned on and off independently.
			\item( ) both lamps go out when one burns out.
			\item( ) the same current flows through all lamps, even if they are different.
			\item(x) it gets brighter if you add more lamps.
			\item( ) both lamps must have the same operating current in amperes.
			\item(x) both lamps must have the same operating voltage in volts.
			\item( ) both lamps must have the same power in watts.
			\item( ) the sum of the operating voltages of the lamps must be equal to the voltage of the power source.
		\end{itemize}
		
		\begin{minipage}{0.45\textwidth}
			\item e) {\bf Series circuit} The circuit is unbranched. When a switch in the series is closed, both lamps turn on. And when it is opened, both lamps turn off.
		\end{minipage}
		\hspace{1 cm}
		\begin{minipage}{0.45\textwidth}
			
			\begin{circuitikz}[american voltages]
				% Battery on the left, two lamps in series at the top, return line at the bottom with a switch
				\draw
				(0,0) to[battery1,l=$V$] (0,3) % Voltage source
				-- (1.5,3) to[lamp,l=$L_1$] (2.5,3) % First bulb
				to[lamp,l=$L_2$] (6,3) % Second bulb
				-- (6,0)
				to[switch,l=$S$] (1,0) % Switch
				-- (0,0); % Return to the negative pole
			\end{circuitikz}
		\end{minipage}
		
		\begin{minipage}{0.45\textwidth}
			{\bf Parallel circuit} The circuit is divided into two branches. A switch can be installed for each lamp. Switch 1 can turn lamp 1 on and off, and switch 2 can turn lamp 2 on and off independently.
		\end{minipage}
		\hspace{1 cm}
		\begin{minipage}{0.45\textwidth}
			\begin{circuitikz}[american voltages]
				% Battery on the left, upper and lower busbars
				\draw
				(0,0) to[battery1,l=$V$] (0,4)
				-- (3.5,4)
				(0,0) -- (3.5,0);
				
				% Left branch: switch + lamp L1 
				\draw 
				(1.5,4) to[lamp,l=$L_1$] (1.5,2) 
				to[switch,l=$S_1$] (1.5,0); 
				
				% Right branch: switch + lamp L2 
				\draw 
				(3.5,4) to[lamp,l=$L_2$] (3.5,2) 
				to[switch,l=$S_2$] (3.5,0); 
			\end{circuitikz} 
		\end{minipage}
		
	\end{itemize}
	
	\begin{center}
		\includegraphics[width=7 cm]{../../viecher/eendcomic.pdf}
		
		Click here to return to the \href{https://www.okuyakl.de/phys/p7escL129/ae129.pdf}{worksheet}
	\end{center}
\end{document}